\section{Discussion}
\label{sec:discussion}

\para{Interpretation.}
The combined evidence supports a clear conclusion: lie-inducing instructions produce a measurable output shift in \gptmodel, and that shift is large enough for practical black-box detection. The effect survives held-out question splits and seed transfer.

\para{Why this matters for deployment.}
Output-only monitoring is easy to integrate in API-based systems. A detector like \methodname can run as a lightweight risk score without requiring proprietary internals, which is useful for auditing, policy enforcement, and post-hoc investigation.

\para{Limitations.}
Our study uses one model family and one primary prompt domain (\truthfulqa). The primary held-out set is moderate in size (72 test outputs), and we do not include human judgments of deception quality. The auxiliary jailbreak-style score is in-sample only and should not be over-interpreted as generalization evidence.

\para{Threats to validity.}
Prompt-template artifacts could contribute to separability, although paired prompts and grouped splitting reduce this risk. Domain concentration in misconception-heavy factual QA may also inflate detectability relative to open-ended dialogue. Finally, adversarial anti-detection prompting was not evaluated.

\para{Broader implications.}
These findings support defense-in-depth: output-level detectors can complement internal probes and policy models. At the same time, detectability does not guarantee prevention; attackers may adapt style to evade fixed detectors, so continual evaluation and recalibration remain necessary.
